\chapter{Стационарный режим при неоднородных граничных условиях}

\section{Постановка индивидуальной задачи}

Рассмотрим стержень длины $l>0$ с теплоизолированной боковой
поверхностью и без источников тепла. Начальная температура равна $x$.
Типы условий теперь задаются индивидуальным вариантом:
первый род при $x=0$ и второй род при $x=l$.

Сохраняем знаки общей постановки с изображения:
\[
 \alpha u(0,t)-\beta u_x(0,t)=\mu(t),\qquad
 \gamma u(l,t)-\delta u_x(l,t)=\nu(t).
\]
Выбираем нормировку
\[
 \alpha=1,\quad\beta=0,\qquad
 \gamma=0,\quad\delta=1,\qquad
 \mu(t)=A,\quad\nu(t)=B.
\]
Следовательно, индивидуальные неоднородные условия имеют вид
\[
\boxed{u(0,t)=A,\qquad -u_x(l,t)=B.}
\]
Знак минус во втором условии обусловлен записью исходного граничного
оператора. Поэтому здесь $B=-u_x(l,t)$, а не $B=u_x(l,t)$.
Параметры $A,B$ --- произвольные вещественные константы.

Полная задача:
\[
\boxed{\begin{cases}
 u_t=a^2u_{xx},&0<x<l,\quad t>0,\\
 u(0,t)=A,\quad -u_x(l,t)=B,&t>0,\\
 u(x,0)=x,&0\le x\le l.
\end{cases}}
\]
Теплоизоляция боковой поверхности означает $b=0$; отсутствие
распределённых источников означает $f=0$. При этом торцы могут
обмениваться теплом с внешней средой: это воздействие задают
неоднородные граничные условия.

\section{Нахождение стационарного профиля}

В стационарном режиме температура не зависит от времени:
\[
 u(x,t)=U(x),\qquad u_t=0.
\]
Следовательно,
\[
 U''(x)=0,\qquad U(0)=A,\qquad -U'(l)=B.
\]
Общее решение дифференциального уравнения имеет вид
\[
 U(x)=C_1x+C_2.
\]
Из левого граничного условия получаем $C_2=A$, а из правого
$-C_1=B$. Поэтому
\[
\boxed{U(x;A,B)=A-Bx,\qquad 0\le x\le l.}
\]
Проверка:
\[
 U''=0,\qquad U(0)=A,\qquad -U'(l)=B.
\]
Стационарный профиль существует и единственен при любых $A,B$.
Действительно, разность двух стационарных решений удовлетворяет
$W''=0$, $W(0)=0$, $W'(l)=0$, откуда $W\equiv0$.

Если условие второго рода нормируют иначе, записывая $u_x(l,t)=B$,
ответ имеет вид $U=A+Bx$. Это другой выбор знака параметра $B$.
В настоящем разделе используется именно форма $-u_x(l,t)=B$
из исходной постановки.

\section{Установление стационарного режима}

Чтобы показать, что найденный профиль является предельным режимом,
выполним замену
\[
 u(x,t)=U(x)+v(x,t)=A-Bx+v(x,t).
\]
Тогда
\[
\begin{cases}
 v_t=a^2v_{xx},&0<x<l,\ t>0,\\
 v(0,t)=0,\quad v_x(l,t)=0,&t>0,\\
 v(x,0)=(1+B)x-A.
\end{cases}
\]
Собственные значения и функции однородной смешанной задачи равны
\[
 k_n=\frac{(n+\tfrac12)\pi}{l},\qquad
 \lambda_n=k_n^2,\qquad X_n(x)=\sin(k_nx),\qquad n=0,1,2,\ldots.
\]
Поэтому решение записывается в виде
\[
 u(x,t)=A-Bx+
 \sum_{n=0}^{\infty}c_n e^{-a^2k_n^2t}\sin(k_nx),
\]
где
\[
 c_n=\frac2l\int_0^l\bigl((1+B)\xi-A\bigr)\sin(k_n\xi)\,d\xi.
\]
Используя $\cos(k_nl)=0$ и $\sin(k_nl)=(-1)^n$, получаем
\[
 \int_0^l\sin(k_n\xi)\,d\xi=\frac1{k_n},\qquad
 \int_0^l\xi\sin(k_n\xi)\,d\xi=\frac{(-1)^n}{k_n^2}.
\]
Таким образом,
\[
\boxed{c_n=\frac2l\left(
 \frac{(1+B)(-1)^n}{k_n^2}-\frac A{k_n}\right).}
\]
Все собственные значения положительны, и переходные моды затухают:
\[
\boxed{u(\cdot,t)\longrightarrow A-Bx\quad(t\to\infty).}
\]
В частности, в норме $L^2(0,l)$
\[
 \|u(\cdot,t)-U\|_{L^2}
 \le e^{-a^2\pi^2t/(4l^2)}\|x-U\|_{L^2}.
\]
Для произвольного начального распределения $\phi\in L^2(0,l)$
формула сохраняется с коэффициентами
\[
 c_n=\frac2l\int_0^l\bigl(\phi(\xi)-A+B\xi\bigr)
 \sin(k_n\xi)\,d\xi.
\]
Следовательно, предельный профиль не зависит от начальных условий.

Начальное распределение $x$ согласовано с обоими граничными условиями
в начальный момент только при $A=0$, $B=-1$.
Для остальных значений решение рассматривается при $t>0$ с
начальным условием в смысле $L^2$; гладкость вплоть до угловых точек
$(0,0)$ и $(l,0)$ не предполагается.
При $A=0$, $B=-1$ начальная функция совпадает со стационарным
профилем, и $u(x,t)=x$ для всех $t$.

\section{Зависимость от параметров}

Профиль линейно зависит от $A$ и $B$:
\[
 U(x;A,B)=A-Bx,\qquad
 \frac{\partial U}{\partial A}=1,\qquad
 \frac{\partial U}{\partial B}=-x.
\]
Изменение $A$ на $\Delta A$ сдвигает весь профиль на одну и ту же
величину. Изменение $B$ на $\Delta B$ меняет температуру на $-x\Delta B$:
влияние отсутствует на левом торце и максимально по модулю на правом.
Точное приращение равно
\[
 \Delta U(x)=\Delta A-x\Delta B,
 \qquad
 \|\Delta U\|_{\infty}\le|\Delta A|+l|\Delta B|.
\]

Температуры концов и средняя температура:
\[
 U(0)=A,\qquad U(l)=A-Bl,\qquad
 \overline U=\frac1l\int_0^lU(x)\,dx=A-\frac{Bl}{2}.
\]
Если $B>0$, температура убывает вдоль стержня;
если $B<0$, возрастает; если $B=0$, профиль постоянен и равен $A$.
Максимальное и минимальное значения равны
\[
 \max U=\max\{A,A-Bl\},\qquad
 \min U=\min\{A,A-Bl\}.
\]
Неотрицательность температуры на всём отрезке эквивалентна условиям
\[
 U(x)\ge0\ \text{для всех }x\in[0,l]
 \quad\Longleftrightarrow\quad A\ge0,\quad A-Bl\ge0.
\]
Строгая положительность на замкнутом отрезке эквивалентна двум
строгим неравенствам. Эти ограничения имеют физический смысл,
если температура понимается как абсолютная; для температуры
относительно выбранного отсчёта отрицательные значения допустимы.

Ненулевой стационарный профиль получается тогда и только тогда, когда
\[
\boxed{U\not\equiv0\quad\Longleftrightarrow\quad (A,B)\ne(0,0).}
\]
Нулевое среднее $A=Bl/2$ не означает тождественного нуля:
при $B\ne0$ профиль меняет знак и остаётся ненулевым.
При $B\ne0$ его нуль расположен в точке $x_*=A/B$,
если эта точка принадлежит отрезку $[0,l]$.

\section{Физическая интерпретация}

Согласно закону Фурье плотность потока в направлении положительной
оси $Ox$ равна
\[
 j_x=-kU'=kB.
\]
Поток через поперечное сечение площадью $S$ равен $Q=kSB$.
Таким образом, $B$ задаёт поток с точностью до множителя $kS$,
а $A$ задаёт температуру левого торца.
При $B>0$ поток направлен вправо; при $B<0$ --- влево.

Внешняя нормаль на левом торце направлена в сторону $-x$,
а на правом --- в сторону $+x$. Поэтому наружные потоки равны
\[
 Q_{\mathrm{out},0}=-kSB,\qquad
 Q_{\mathrm{out},l}=kSB,
 \qquad Q_{\mathrm{out},0}+Q_{\mathrm{out},l}=0.
\]
В стационарном режиме тепло не накапливается: поток, поступающий
через один торец, выходит через другой. Ненулевой градиент температуры
и ненулевой поток совместимы со стационарностью.
Коэффициент $a^2$ не входит в формулу $U=A-Bx$, но определяет скорость
затухания переходного процесса.
